% Introduction for Chapter 1; kept in main.tex.
\newcommand{\ChapterOneIntroduction}{%
\section*{Introduction}
Our aim in this chapter is to develop the familiar theories of analytic functions of one complex variable and Riemann surfaces in a way that generalises well to the several variable theory. In \hyperref[sec:1:3]{\S3} we see how the existence theory of the Cauchy--Riemann equations can be used to prove the Mittag-Leffler theorem and also how the topology of domains in $\mathbb C$ naturally enters into the proof of the Weierstrass theorem. In \S\S4,5 we show how the theory of holomorphic line bundles may be used to reformulate some of the classical problems in Riemann surface theory. We also define the Cauchy--Riemann equations on an arbitrary Riemann surface and indicate how they are related to the problem of constructing meromorphic functions with specified divisors. In an appendix we prove a number of classical results, including the Runge approximation theorem. We use the Runge theorem to construct solutions of the Cauchy--Riemann equations.
}

% Introduction for Chapter 2; kept in main.tex.
\newcommand{\ChapterTwoIntroduction}{%
\section*{Introduction}
In this Chapter we develop the basic theory of analytic functions of several complex variables with particular reference to the problem of extension of analytic functions. Section 1 is a straightforward generalisation of the one variable theory and contains such theorems as the power series representation of an analytic function. In section 2 we prove a Riemann removable singularities theorem for analytic functions of more than one variable. The material in section 3 is in sharp contrast to the one variable theory and includes the theorem of Hartog's that implies, for example, that every analytic function on the punctured Euclidean disc in $\mathbb C^n$, $n>1$, extends analytically to the whole disc. In section 4 we define domains of holomorphy and prove their equivalence with holomorphically convex domains. Section 5 is devoted to various pseudoconvexity properties that a domain of holomorphy possesses. In section 6 we discuss the Bergman kernel function of a domain and in section 7 we make a preliminary study of the Cousin problems.
}

% Introduction for Chapter 3; kept in main.tex.
\newcommand{\ChapterThreeIntroduction}{%
\section*{Introduction}
In this chapter we shall make a study of algebraic and analytic properties of the ring $\mathbb C\{z_1-a_1,\ldots,z_n-a_n\}$ of convergent power series at a point $(a_1,\ldots,a_n)\in\mathbb C^n$. It turns out that there is a rich and deep relationship between algebraic properties of power series rings and local properties of analytic functions and analytic sets. This relationship is expressed in its most elegant and powerful form using the concept of ``coherence'' which we introduce later in Chapter~\ref{ch:7}. In this chapter we show how to define meromorphic functions of more than one complex variable. We also prove some elementary facts about local properties of analytic sets and some not so elementary results about modules over $\mathbb C\{z_1-a_1,\ldots,z_n-a_n\}$. The main technical result we need is the most important Weierstrass Preparation theorem which we prove in section 2.
}

% Introduction for Chapter 4; kept in main.tex.
\newcommand{\ChapterFourIntroduction}{%
\section*{Introduction}
Much of this chapter is taken up with the discussion of specific examples of complex manifolds. After general definitions given in section 1, we consider complex submanifolds of $\mathbb C^n$ in section 2. We prove that the open Euclidean disc and polydisc are biholomorphically inequivalent. Included in this section are also definitions of the classical domains and Stein manifolds. In section 3 we consider projective algebraic manifolds and in section 5 complex manifolds defined as the quotient of a properly discontinuous group action. In section 4 we define complex tori and prove that every one dimensional complex torus is biholomorphic to a cubic curve. In section 6 we develop the structure theory of complex hypersurfaces and in section 7 we define blowing up and give some indication of its use in desingularising analytic sets and in the classification theory of complex manifolds.
}


% The chapter introductions are intentionally kept in main.tex.
\newcommand{\ChapterFiveIntroduction}{%
% Introduction from PDF page 10; printed page 1
\paragraph*{Introduction}
In sections 1 and 2 we cover basic linear algebra and calculus on differential manifolds. We define the complexification of a real vector space in section 3 and in section 4 we develop the main results of complex linear algebra that we need in the sequel. After giving some general facts about complex and holomorphic vector bundles in section 5 and constructing the tangent and cotangent bundles of a complex manifold in section 6 we reach the heart of the chapter in section 7 with the construction of the $\bar\partial$-operator on an arbitrary complex manifold. In section 8 we prove the Dolbeault--Grothendieck Lemma and, with a little more effort, deduce that the Cousin I and II problems are always solvable on polydiscs. Section 9 is devoted to the discussion of a number of important examples on compact complex manifolds. Thus we construct the Euler sequence for projective space, prove some basic results about linear systems and their relationship with holomorphic line bundles and conclude with a discussion of theta functions and complex tori. Finally, in section 10, we define the concepts of $q$-pseudoconvexivity, $q$-completeness and weakly positive vector bundles.

}
\newcommand{\ChapterSixIntroduction}{%
\section*{Introduction}
In Section 1 of this chapter we present the basic definitions and constructions of sheaf theory with many motivating examples. In section 2 we give an application of sheaf theory to prove the existence and uniqueness of the envelope of holomorphy of a Riemann domain. In section 3 we define the sheaf cohomology groups of a sheaf of groups over a paracompact space using fine resolutions. Amongst the most important results we prove are Leray's theorem and the existence of a canonical, natural isomorphism between \v{C}ech cohomology and sheaf cohomology. We conclude with a number of important examples and computations involving the 1st. Chern class.

}
\newcommand{\ChapterSevenIntroduction}{%
% Introduction from PDF page 136; printed page 127
\paragraph*{Introduction}
In section 1 we define coherence for sheaves of $\mathcal{O}$-modules and prove Oka's theorem. In section 2 we prove Cartan's theorems A and B for Stein manifolds assuming exactness of the $\bar\partial$-sequence. The remainder of the Chapter is devoted to applications of Theorems A and B. We prove Cartan and Serre's finiteness theorem in section 3 and Grauert's finiteness theorem for coherent sheaves on s.l.p. domains in section 4. Using the finiteness theorem of Cartan-Serre, we prove Serre's theorems A and B for coherent sheaves on complex projective space in section 5. We give a number of applications including Grothendieck's theorem on the splitting of holomorphic vector bundles on the Riemann sphere. Finally in section 6 we prove Kodaira's embedding theorem, following Grauert, and conclude by showing that complex tori that admit a Riemann form are algebraic.

}

